\documentclass{amsart}

\usepackage[margin=1in]{geometry}
\usepackage[width=\linewidth]{caption}
\usepackage{tikz,tikz-cd}
\usetikzlibrary{calc}


\definecolor{darkblue}{rgb}{0,0,0.4} 
 \usepackage[colorlinks=true, citecolor=darkblue, filecolor=darkblue, linkcolor=darkblue,urlcolor=darkblue]{hyperref} 
\usepackage[all]{hypcap}
\usepackage{xr-hyper}

\usepackage{ifthen,calc,amsmath,amssymb,mathtools,enumitem}

\usepackage{xcolor,amssymb}
\usepackage{wasysym}
\usepackage{pdflscape}

\usepackage{pdfpages}

\newtheorem*{problem*}{Problem}

\newcommand{\from}{\colon}
\newcommand{\FF}{\mathbb{F}}
\newcommand{\ZZ}{\mathbb{Z}}
\newcommand{\RR}{\mathbb{R}}
\newcommand{\defeq}{\coloneqq}
\newcommand{\bdy}{\partial}

\begin{document}


\begin{problem*}
  Let $G=\langle a_1,\dots,a_r\rangle$ be a free group. A word on the
  letters $L=\{a_1,a_1^{-1},\dots,a_r,a_r^{-1}\}$ is called reducible
  if it equals $1\in G$. A word is called quasi-reduced if it does not
  contain three consecutive letters $\ell\ell^{-1}\ell$ for any
  $\ell \in L$.

  Fix a reducible quasi-reduced word $w=w_1\cdots w_{2n}$ (with $w_i\in L$). A
  matching $M$ consists of $n$ intervals, each properly embedded in
  $\RR\times\RR_+$, satisfying the following:
  \begin{itemize}[leftmargin=*]
  \item Any two intervals intersect
    transversally, and there are no triple points.
  \item $\bdy M=\{1,\dots,2n\}\times 0\subset\RR\times 0\subset
    \RR\times\RR_+$.
  \item For each component $R$ of
    $\RR\times\RR_+\setminus M$, and for each component $A$ of $M\cap \bdy\overline{R}$, $\bdy A$ consists of some two points $(i,0)$ and $(j,0)$ with $w_i=w_j^{-1}$. 
  \end{itemize}
  A $k$-crossing
  matching is a matching where the total number of self-intersections
  is $k$.

  Let $F_s$ be the following CW complex whose $k$-cells correspond to
  isotopy classes of $k$-crossing matchings. Given a $k$-crossing
  matching $M$ with crossing set $X(M)$, the associated $k$-cell is
  given by $C(M)\defeq \prod_{c\in X(M)}I_c$, where $I_c$ is an
  interval whose endpoints correspond to the two ways of locally
  resolving $M$ near $c$. The attaching maps are given as follows: On
  a facet of the form
  \[
    a\times\prod_{d\in X(M)\setminus\{c\}}I_{d},\qquad c\in X(M),a\in \bdy I_c,
  \]
  let $M_a$ be the $(k-1)$-crossing matching obtained by resolving $M$
  at $c$ according to $a$. Then the $(k-1)$-cell associated to $M_a$
  is given by
  $C(M_a)=\prod_{d\in X(M_a)}I_d=\prod_{d\in X(M)\setminus\{c\}}I_{d}$, and the
  above facet maps to this $(k-1)$-cell canonically.

  Is $F_s$ contractible?
  \end{problem*}

\begin{proof}[History]
  \renewcommand{\qedsymbol}{} A more general version of this complex
  is being studied by Robert Lipshitz, Lenny Ng, and Sucharit Sarkar
  for their upcoming paper ``A stable homotopy refinement of
  Legendrian contact homology'', scheduled to appear in mid-2026. They
  have proved that $F_s$ is always $1$-connected, which is needed for
  some applications in their paper. If it were contractible, that
  would have strengthened some of the results in their paper, but much
  to their disappointment, they discovered that $F_s$ is not even
  $2$-connected in general.
\end{proof}

\begin{proof}[Solution]
  \begin{figure}[h]
    \centering
    \begin{tikzpicture}[xscale=4,yscale=-2.5]
      \foreach\a/\b [count=\c from 0] in{0/0,0/2,2/0,2/2,1/1,0.5/1.5,1.5/0.5,2.5/1}{

        \node[anchor=south] (v\a\b) at (\a,\b) {        
          \begin{tikzpicture}[xscale=0.3,yscale=-0.5]

            \foreach \i in {0,...,5}{
              \coordinate (l\i) at (\i,0);
              \coordinate (a\i) at (\i,-0.5);
            }


            \ifnum\c=0
            \draw[thick] (a0) to[out=-90,in=-90] (a1);
            \draw[thick] (a2) to[out=-90,in=-90] (a3);
            \draw[thick] (a4) to[out=-90,in=-90] (a5);
            \fi
            \ifnum\c=1
            \draw[thick] (a0) to[out=-90,in=-90] (a1);
            \draw[thick] (a2) to[out=-90,in=-90] (a5);
            \draw[thick] (a4) to[out=-90,in=-90] (a3);
            \fi
            \ifnum\c=2
            \draw[thick] (a0) to[out=-90,in=-90] (a3);
            \draw[thick] (a2) to[out=-90,in=-90] (a1);
            \draw[thick] (a4) to[out=-90,in=-90] (a5);
            \fi
            \ifnum\c=3
            \draw[thick] (a0) to[out=-90,in=-90] (a5);
            \draw[thick] (a2) to[out=-90,in=-90] (a1);
            \draw[thick] (a4) to[out=-90,in=-90] (a3);
            \fi
            \ifnum\c=4
            \draw[thick] (a0) to[out=-90,in=-90] (a5);
            \draw[thick] (a2) to[out=-90,in=-90] (a3);
            \draw[thick] (a4) to[out=-90,in=-90] (a1);
            \fi
            \ifnum\c=5
            \draw[thick] (a0) to[out=-90,in=-90] (a3);
            \draw[thick] (a1) to[out=-90,in=-90] (a5);
            \draw[thick] (a2) to[out=-90,in=-90] (a4);
            \fi
            \ifnum\c=6
            \draw[thick] (a0) to[out=-90,in=-90] (a4);
            \draw[thick] (a1) to[out=-90,in=-90] (a3);
            \draw[thick] (a2) to[out=-90,in=-90] (a5);
            \fi
            \ifnum\c=7
            \draw[thick] (a0) to[out=-90,in=-90] (a2);
            \draw[thick] (a1) to[out=-90,in=-90] (a4);
            \draw[thick] (a3) to[out=-90,in=-90] (a5);
            \fi
            
          \end{tikzpicture}};

      }

      \draw[ultra thick] (v00)--(v02) (v22)--(v20);
      \draw[ultra thick] (v00)--(v00-|v20.west) (v02)--(v02-|v22.west);
      \draw[ultra thick] (v00)--($(v00)!0.35!(v22)$) ($(v00)!0.6!(v22)$)--(v22);
      
    \end{tikzpicture}
    \caption{The space $F_{aa^{-1}bb^{-1}aa^{-1}cc^{-1}aa^{-1}}$. We
      have drawn the five $0$-crossing matchings for the $0$-cells and
      the three $2$-crossing matchings for the $2$-cells, but not the
      six $1$-crossing matchings for the
      $1$-cells.}\label{fig:menorah}
  \end{figure}

  
  The answer is no. The space $F_s$ is not $2$-connected in general. For instance, if $r\geq 3$ and $a=a_1,b=a_2,c=a_3$, and $w=aa^{-1}bb^{-1}aa^{-1}cc^{-1}aa^{-1}$, $F_w$ consists of five $0$-cells, six $1$-cells, three $2$-cells,
  and is homeomorphic to $S^2$, as illustrated in Figure~\ref{fig:menorah}.
\end{proof}



\end{document}

